% 第7章 结论与展望
% 红线：只写证据支撑的结论；计划/待补事项明确标"待完成"，不写"显著提高"等无证据表述
% 证据：前六章各节注释；素材总览见 materials/素材索引.md
\section{结论与展望}

\subsection{主要结论}

本文围绕《AI辅助下的DRAM存储器单粒子效应评估研究》，建立了从 Geant4 粒子输运到 TCAD 器件响应的可审计联合建模与数据链路，统一了指标与标签口径，构建了带 run ID 与哈希的冻结数据集，并设计了 R$\rightarrow$A$\rightarrow$F 数据驱动预测链路与物理约束消融、分组交叉验证、冻结盲测协议。已有证据支持以下结论：

\begin{enumerate}
  \item 在 v3 冻结设计（未校准二维代表模型、固定 5\,MeV、两角度、固定入射位置）下，64 个独立事件 / 128 条配对响应可按匹配对照与 $\varepsilon=5\mV$ 主阈值严格分类（96 正确 / 20 判据误读 / 8 临界 / 4 数值不确定），预登记数值细化在容差内、失败案例（3\ns 预充偏差 13.297\uV）被保留而非掩盖——方法链路的可复现性得到验证。
  \item 基于沉积电荷、栅长与角度特征的岭回归基线在八工况留一交叉验证中显著优于均值基准（$\Delta Q$ 折 RMSE 0.103/0.113\fC，$\Delta M$ 0.955/1.012\mV），域内 $\Delta M$ MAE 约 0.53--0.55\mV，支持"给定源项与基线条件下，低维回归可作 TCAD 响应代理"的域内可行性。
  \item v4/r0.6 阶段在厚硅输运与三维源接口上取得阶段性资格证据（212 初级事件、9 条核非弹性记录、三个三维源账本闭合、细网格 A0/A1 全帧通过），P0 四项资格门与 D0--D3 第一波正在推进。
\end{enumerate}

\subsection{局限与待补}

本文的器件结论限于未校准代表模型，不代表真实 DDR5；判据误读不等于持久位翻转；配对响应与复跑不增加独立样本。以下事项\textbf{待完成}，不在本稿中宣称：v3 插值/外推留出的全量独立验收与新模型的冻结盲测；r0.6 合格数据集导出与 R/A/F 正式训练；物理约束消融的定量比较；带注量/分母的截面与错误率；与导师课题组 DDR5 实验原始数据（M01/M02）的外部对比；真实器件 I--V/保持/读出校准（M04/M05）。

\subsection{后续工作}

后续按 r0.6 排批推进 P0 资格门与 D0--D3/LHS 合格数据，导出可训练快照后启动 R/A/F 训练与消融、分组 CV 及冻结盲测；同步推进向导师交接实验原始数据、器件校准依据与毕业设计模板（M01/M02/M04/M06），补齐统计分母后开展截面/错误率与外部验证，并按学校模板完成三门课成稿。
